%! Justifying a Claim Based on a Confidence Interval for a Population Proportion — Unit 6, Lesson 6.3 — Warm-Up
\documentclass[10pt]{article}
\usepackage{apstats-article}
\usepackage{apstats-boxes}
\newcommand{\TallMath}[1]{$\displaystyle #1\rule[-1.4em]{0pt}{3.2em}$}

\begin{document}

\pageheader{Unit 6, Lesson 6.3}{Warm-Up}
\namedateperiod

\vspace{0.15in}

A random sample of 360 college students found 198 who said they regularly eat breakfast.

\begin{enumerate}
  \item Verify the three conditions for a one-sample $z$-interval.
        \begin{itemize}
          \item Random: \blank{7cm}
          \item Large Counts: \blank{7cm}
          \item 10\% rule: \blank{5cm}
        \end{itemize}

  \item Construct a 95\% confidence interval for $p$, the true proportion of all college
        students who regularly eat breakfast. Show all work.

        $\hat p = $ \blank{2.5cm}

        \[
          \hat p \pm z^* \sqrt{\frac{\hat p(1-\hat p)}{n}}
          = \blank{1.5cm} \pm \blank{1.5cm}\sqrt{\frac{\blank{2cm}}{\blank{1.5cm}}}
          = (\blank{2.5cm},\; \blank{2.5cm})
        \]

  \item Based on your interval, is the value $p = 0.60$ plausible? Is $p = 0.52$ plausible?
        Explain briefly.

        \writelines{3}
\end{enumerate}

\end{document}
